\subsection{Sensitivity to the accepted input information}\label{sec:sourcepolicy}
Point-withdrawal tests keep the primary data, weights, supplements and 48
non-point bounds fixed, replacing selected point utilities by $[0,1]$.
Of 112 points, 64 are dated 2026, two 2025 and 46 earlier. The 28
non-\texttt{ok} points comprise 16 \texttt{MODELLED}, eight \texttt{STALE},
three \texttt{SUSPECT} and one \texttt{NATIONAL}. Groups cover dates before
2020, after 2025, non-\texttt{ok} flags and their union; overlaps count once.

\begin{table}[!ht]\centering
\caption{Recommendation after withdrawing selected submitted point information. Width is the mean baseline-weight score range on the 0--100 scale.}\label{tab:vintage}
\begin{tabular}{l r l l r}\toprule
Withdrawn group & Cells & Baseline & Equal & Width\\\midrule
Retained input & 0 & Zurich & Zurich & 16.05 \\
Before 2020 & 13 & Singapore & Zurich & 21.12 \\
After 2025 & 64 & Zurich & Zurich & 48.41 \\
Non-ok flags & 28 & Berlin & Paris & 29.11 \\
Union & 95 & Singapore & Singapore & 68.78 \\
\bottomrule\end{tabular}\end{table}


Withdrawing 28 flagged points selects Berlin under baseline weights and Paris
under equal weights. Removing 13 pre-2020 points selects Singapore under
baseline weights and retains Zurich under equal weights. Removing 64
later-dated points retains Zurich under both, widening the baseline mean
range from 16.05 to 48.41. The 95-point union selects Singapore under both.
All sixteen remain possible champions. Labels and cutoffs are not error
probabilities or proof that flagged entries are wrong.

Withholding Zurich's homicide-rate point selects Singapore under baseline
weights and Oslo under equal weights; withholding Amsterdam's retains Zurich
under both. Other entries remain fixed, without a replacement rate.
Varying only Zurich's point rate gives regret ties at about 0.23002 per 100,000
under baseline weights and 0.90473 under equal weights. Below each respective
cutoff Zurich uniquely minimizes regret; above it Zurich does not. These are
hypothetical cutoffs, not measured rates or confidence limits.

Separately withdraw commute proxies by replacing their travel intervals with
$[0,1]$, fixing other entries. Withholding Zurich's alone, both identified
proxies, or all twelve supplied commute means selects Singapore under baseline
weights; equal weights retain Zurich. Withholding Dublin's alone retains Zurich
under both. All sixteen remain possible champions. The baseline recommendation
therefore depends on retaining Zurich's work-commute proxy.
